% Gauss, Werke X/1 (1917), native65–66, printed58–59.
\SourceLeaf{065}\NoSourceHead
\SourceGroup{X}
\HistoricTitle{ZUR THEORIE DES BIQUADRATISCHEN RESTES $1+i$.}
\HistoricCitation{[Zwei Zettel in Ec3, Kapsel 43.]}
\NumberedItem{1}
\par\indent Modulus $a+bi=p$, $b$ gerade und positiv, $aa+bb=D$.

1. Es sei $\Omega$ der Inbegriff der Zahlen $x+iy$, wo
\[
\frac{x+iy}{a+bi}=\xi+i\eta,
\]
\noindent und $\xi,\eta$ zwischen $0$ und $\tfrac12$.

2. Man theile $\Omega$ in zwei Classen $M$ und $N$, wo für $M$ $\xi-\eta$ positiv, für $N$ negativ ist, und nenne die Anzahl der in $M$ und $N$ enthaltenen Zahlen $m,n$. Man hat dann
\[
m+n=\tfrac14(D-1).
\]
\par\noindent\hspace*{\parindent}3. Der Decident von
\[
\begin{array}{rccl}
1+i&\text{ist dann}&\equiv&3m+2n\\
1-i&&\equiv&2m+n\\
-1+i&&\equiv&3n\\
-1-i&&\equiv&m
\end{array}
\]
\par\noindent\hspace*{\parindent}4. Durch Induction haben wir gefunden, dass wenn unter den Zahlen
\[
\left[\frac ab\right],\left[\frac{3a}b\right],\left[\frac{5a}b\right],\left[\frac{7a}b\right],\ldots,\left[\frac{(b-1)a}b\right]
\]
\clearpage

\SourceLeaf{066}\ArithmeticHead{59}{R}{ZUR THEORIE DES BIQUADRATISCHEN RESTES $1+i$.}
\noindent zusammen $\mu$ gerade und $\nu$ ungerade sind, man
\[
m-\mu=n-\nu
\]
\noindent haben werde, wodurch $m$ und $n$ bestimmt sind.

5. Beispiel $p=9+14i$
\[
\left[\frac9{14}\right],\left[\frac{27}{14}\right],\left[\frac{45}{14}\right],\left[\frac{63}{14}\right],\left[\frac{81}{14}\right],\left[\frac{99}{14}\right],\left[\frac{117}{14}\right]=0,\ 1,\ 3,\ 4,\ 5,\ 7,\ 8;
\]
\[
\mu=3,\ \nu=4
\]
\noindent\hspace*{\parindent}[Beispiel] $p=5+12i$
\[
\left[\frac5{12}\right],\left[\frac{15}{12}\right],\left[\frac{25}{12}\right],\left[\frac{35}{12}\right],\left[\frac{45}{12}\right],\left[\frac{55}{12}\right]=0,\ 1,\ 2,\ 2,\ 3,\ 4;
\]
\[
\mu=4,\ \nu=2.
\]
\par\noindent\hspace*{\parindent}6. Man hat $\mu+\nu=\tfrac12b$,
\[
\nu=\left[\frac ab\right]+\left[\frac{3a}b\right]+\left[\frac{5a}b\right]+\ldots-2\left[\frac a{2b}\right]-2\left[\frac{3a}{2b}\right]-2\left[\frac{5a}{2b}\right]\ldots
\]
\noindent Der erste Theil wird $=\tfrac14(a-1)b$, insofern $a$ positiv,
\[
\begin{aligned}
\nu={}&\tfrac14(a-1)b+2\left[\frac ab\right]+2\left[\frac{2a}b\right]+2\left[\frac{3a}b\right]+\ldots+2\left[\frac{\tfrac12b\mathpunct{.}a}b\right]\\
&-2\left[\frac a{2b}\right]-2\left[\frac{2a}{2b}\right]-2\left[\frac{3a}{2b}\right]-\ldots-2\left[\frac{ba}{2b}\right]\\
={}&\tfrac14(a-1)b+2\varphi(b,a)-2\varphi(2b,a)\\
={}&-\tfrac14(a-1)b-2\varphi(a,b)+2\varphi(a,2b).
\end{aligned}
\]
\par\noindent\hspace*{\parindent}Nun ist
\[
\begin{aligned}
2\varphi(a,2b)-2\varphi(a,b)&=\tfrac12(a-1)(b-1)\\
&\quad-4\left\{\left[\frac ba\right]+\left[\frac{3b}a\right]+\left[\frac{5b}a\right]+\ldots+\left[\frac{\tfrac12(a-3)b}a\right]\right\}
\end{aligned}
\]
\centerline{wenn $a\equiv1\ (\mathrm{mod.}\ 4)$,}
\[
\begin{aligned}
{}={}&\tfrac12(a+1)(b-1)\\
&-4\left\{\left[\frac ba\right]+\left[\frac{3b}a\right]+\left[\frac{5b}a\right]+\ldots+\left[\frac{\tfrac12(a-1)b}a\right]\right\}
\end{aligned}
\]
\centerline{wenn $a\equiv3\ (\mathrm{mod.}\ 4)$,}
\noindent also
\[
\begin{array}{cll}
\equiv0\ \mathrm{mod.}\ 4,&\text{wenn $a$ von der Form}&8n+1,8n+7\\
\equiv2\ \mathrm{mod.}\ 4,&\text{wenn $a$ von der Form}&8n+3,8n+5.
\end{array}
\]
\par\vfill\noindent\hfill8*\par\clearpage
